\paragraph{\trFibrationWithStriplikeEnds。}\label{sub:fibrations-striplike-ends}
设 $M$ 是一个\trExactSymplectic\trCorneredBoundary\trManifold，而 $\hat\pi\colon\hat E\to\hat S$ 是一个\trLefschetzFibration。设在基上有一有限的\trPointedBoundary{}集合 $\Sigma=\Sigma^-\sqcup\Sigma^+$，并且在每个这样的点处，都有一个把\trFiber{}$\hat E_\zeta$ 等同到 $M$ 的等同。我们还将假定，该\trFibration{}在诸\trPointedBoundary{}附近是局部平凡的。这意味着每个 $\zeta$ 都有一个邻域 $\hat U=\hat U_\zeta\subset\hat S$ 使得
\begin{equation*}
(\hat\pi|\hat U\colon\hat E|\hat U\to\hat U,\phi_{\hat E},\theta_{\hat E},I_{\hat E})
=(\hat U\times M\to\hat U,\phi_{\hat S}+\phi_M,\theta_{\hat S}+\theta_M,I_{\hat S}\times I_M).
\end{equation*}
照常取 $S=\hat S\setminus\Sigma$；并对每个 $\zeta$ 选取一个\trStriplikeEnd{}$\epsilon_\zeta$，其像落在 $U_\zeta=\hat U_\zeta\setminus\{\zeta\}$ 之内。于是，$\hat\pi$ 的限制 $\pi\colon E\to S$ 便在诸\trStriplikeEnd{}之上自带\trCanonical{}的\trTrivialization{}$Z^\pm\times M\to E$。我们将把这类\trObject{}称作\emph{\trLefschetz\trFibrationWithStriplikeEnds}（Lefschetz fibrations with strip-like ends）。
\begin{remark}\label{thm:local-triviality-deformation}
在许多情形，\trPointedBoundary{}附近的局部平凡性并不\trApriori{}成立，但可以按如下方式用一个\trDeformation{}来补救。设 $\hat\pi\colon\hat E\to\hat S$ 是一个在 $\partial^h\hat E$ 附近平凡的\trLefschetzFibration。任取某个 $\zeta\in\partial\hat S$，并令 $M=\hat E_\zeta$。此时，通过改动这根\trFiber{}附近的\trSymplectic{}结构与\trAlmostComplex{}结构，便可在 $\zeta$ 附近达到局部平凡（更确切地说，按\cref{sub:isomorphisms-deformations}的术语，这将是一个逐\trFiber{}常值的\trDeformation）。这类似于\cref{thm:transport-isotopy-realization}：在一个局部\trTrivialization{}$\hat U\times M\to\hat E|\hat U$ 中，我们有
\begin{equation}\label{eq:theta-local-trivialization}
\theta_{\hat E}=\theta_M+\theta_{\hat S}+\kappa+dR,
\end{equation}
其中 $\kappa$ 与函数 $R$ 都在 $\hat U\times\partial M$ 附近为零。我们想把 \eqref{eq:theta-local-trivialization} 中的这两项换成 $\psi(\kappa+dR)$，其中 $\psi$ 是 $\hat U$ 上一个在紧邻 $\zeta$ 处为零的适当截断函数。修改后的 1-形式的外\trDifferential{}在\trTotalSpace{}上可能不再是\trSymplectic{}的，但这很容易补救，只需再加上基上某个适当项的\trPullback{}即可。同样地，若在局部\trTrivialization{}中考察\trAlmostComplex{}结构{}$I_{\hat E}$，它在 $\hat U\times\partial M$ 附近变得等于 $I_{\hat S}\times I_M$，而且可以容易地对其加以修改，使它在整根\trFiber{}$\hat E_\zeta$ 附近都具有同样的乘积形式。
\end{remark}
我们需要把\trLagrange\trBoundaryCondition{}与\trRelative\trPerturbationDatum{}这两个概念扩展到\trStriplikeEnd{}的语境。对其中第一个，取一个\trProjection{}到基 $S$ 上为\trProper{}的且满足 \eqref{eq:lagrangian-boundary-condition} 的 $F\subset E$。此时，\trFamily{}$F_z$ 自动在每个\trStriplikeEnd{}上\trLocallyConstant，这意味着存在某一对\trLagrange\trSubmanifold{}$(L_{\zeta,0},L_{\zeta,1})$ 使得 $F_{\epsilon_\zeta(s,k)}=L_{\zeta,k}$ 对一切 $s$ 成立。此外，此时必然有 $\alpha_F$ 在诸\trStriplikeEnd{}上为零。现在设对每一对 $(L_{\zeta,0},L_{\zeta,1})$ 都已取定一个\trFloerDatum{}$(H_\zeta,J_\zeta)$。此时，\trRelative\trPerturbationDatum{}便是照旧的一对 $(K,J)$，但附加一条限制：在每个\trStriplikeEnd{}上，$K(s,t,x)=H_\zeta(t,x)\,dt$ 而 $J(s,t,x)=i\times J_\zeta(t,x)$。取 \eqref{eq:section-equation} 的一个解 $u$，并利用每个端上的\trTrivialization{}把它写成 $(u\circ\epsilon_\zeta)(s,t)=(s,t,u_\zeta(s,t))$，其中 $u_\zeta\colon Z^\pm\to M$。此时 \eqref{eq:section-equation} 便化归为 $u_\zeta$ 的\trFloerEquation，因而对 $u_\zeta$ 施加向某个 $y_\zeta\in\scr{C}(L_{\zeta,0},L_{\zeta,1})$\trAsymptotic\trConvergence{}的条件就有意义了。我们把所得的\trModuliSpace{}记作 $\scr{M}_{E/S}(\{y_\zeta\})$。假定适当的\trRegularity{}假设得到满足，便可以定义一个链\trHomomorphism
\begin{equation}\label{eq:cphi-fibration-map}
C\Phi_{E/S}\colon
\bigotimes_{\zeta+\in\Sigma+}CF^{\mathrm{pr}}(L_{\zeta+,0},L_{\zeta+,1})
\longrightarrow
\bigotimes_{\zeta-\in\Sigma-}CF^{\mathrm{pr}}(L_{\zeta-,0},L_{\zeta-,1})
\end{equation}
一如 \eqref{eq:cphi-map}，其方式是对 $\scr{M}_{E/S}(\{y_\zeta\})$ 中的\trIsolated{}点作 mod 2 计数。在\trFloerCohomology{}上诱导出的映射 $\Phi_{E/S}$ 与\trPerturbationDatum{}的选取无关，并且服从 TQFT 公理的一个直截了当的推广。特别地，其中有一条\trCompositionLaw：把两个\trLefschetz\trFibration{}$\pi_k\colon E_k\to S_k$ \trGluing{}在一起，办法是把 $S_1$ 的一个出向端与 $S_2$ 的一个入向端认同起来（当然，诸 $S_k$ 上的\trSymplecticStructure{}通常不会彼此一致，但这可以像\cref{thm:local-triviality-deformation}中那样通过局部拼合来处理）。
这一框架同时包含了先前考察过的 $S$\trCompact{}的情形与 §8 的理论（后者相当于只允许平凡的\trFibration{}$E=S\times M$）。
通过借鉴这两种情形中的想法，可以把它再朝几个方向加以扩展：
\begin{itemize}
\item \emph{\trMovingBoundaryCondition。}与\trCompact{}情形中一样，只是附加一条限制：诸 $F_z$ 需要在每个\trStriplikeEnd{}上关于 $z$ \trLocallyConstant。
\item \emph{\trParametrization{}的\trModuliSpace}（parametrized moduli spaces）。\trLefschetzFibration{}的\trFamily{}及其相伴的\trModuliSpace{}已经出现在基本构造之中，人们用它们来为\trRelative\trPerturbationDatum{}的不同选取所对应的诸映射 \eqref{eq:cphi-fibration-map} 之间提供\trChainHomotopy。其一般结构是：考察一个\trFamily{}$E^r\to S^r$，带有\trLagrange\trBoundaryCondition{}$F^r$ 与依赖于某个参数空间 $r\in R$ 的\trPerturbationDatum{}$(K^r,J^r)$，其中\trStriplikeEnd{}的结构关于 $r$ 是常值的。为简单起见，设 $R$ \trCompact。\trFloerCochainGroup{}之间相应的映射将满足形状如下的方程：$\partial\circ(\text{映射})+(\text{映射})\circ\partial+(\text{边界项})=0$，其中边界项来自把我们的\trFamily{}限制到 $\partial R$ 上。特别地，当 $R=[0;1]$ 时，我们得到两个映射之间的一个\trChainHomotopy，分别对应于两个端点。稍后我们将见到若干双参数的例子（看得更远一点，可以沿 §9 的路线建立一套一般理论，但所得代数结构的含义只有部分得到了理解）。
\item \emph{\trEvaluationMap。}除\trStriplikeEnd{}之外，还可以在 $\partial S$ 上取定一个有限点集 $\Pi=\Pi^-\sqcup\Pi^+$，并考虑在这些点处的\trEvaluationMap{}$\mathrm{ev}_{z_\pm}(u)=u(z_\pm)$。这就给出了一个\trInvariant，它如 \eqref{eq:evaluation-tqft-map} 那样，在\trTensorProduct{}中含有附加的因子 $H^*(F_{z_\pm};\KK)$。例如，在\cref{thm:local-model-sections}的局部模型中，任一点处的\trEvaluationMap{}都可以认同为从切球丛到基的\trProjection{}$S(TS^n)\to S^n$。从根本上说，引入\trMarkedPoint{}并不会使理论更为一般，因为（撇开技术性细节不谈，并使用 PSS \trIsomorphism）它们可以换成\trStriplikeEnd；不过这在计算上往往更为方便。对这种情形，\trCompositionLaw{}还有一个变体，即沿着\trMarkedPoint{}取基\trSurface{}的\trBoundaryConnectedSum。
\end{itemize}
